\section{Introduction}
\label{sec:introduction}

A retinal vessel segmentation model that averages 0.78 Dice across a test
set tells a clinician very little about any individual image. Some images
are segmented near-perfectly; others contain errors that, if acted upon,
could affect screening decisions for diabetic retinopathy or glaucoma. The
aggregate metric hides this variance, and a calibrated probability map does
not, on its own, tell the clinician what to \emph{do} about it. What a
clinical deployment actually needs is a \emph{policy}: for each pixel,
should the prediction be trusted, flagged for review, or handed off to a
human?

\paragraph{Estimation has outpaced decision.} Research on uncertainty in
medical image segmentation has concentrated almost entirely on
\emph{estimation}: better posteriors via Bayesian approximations
\citep{gal2016dropout}, ensembles \citep{lakshminarayanan2017simple}, and
test-time augmentation \citep{wang2019aleatoric}, scored by calibration
metrics such as ECE and NLL. These lines of work ask \emph{how well} we
estimate uncertainty; they rarely ask what happens next. Selective
prediction \citep{geifman2017selective,el2010foundations} formalizes
abstention in classification, but its application to dense pixel-level
prediction---and its coupling with modern uncertainty estimators---remains
underdeveloped. The result is a literature in which uncertainty maps are
produced, tabulated, and then left unused. The missing layer is
\emph{decision}: a policy that converts an uncertainty map into an action.
\textbf{Uncertainty without decisions is evaluation without consequence.}

\paragraph{A decision-aware perspective on uncertainty.} We reframe the
problem. The key question is not how well uncertainty is estimated, but
whether it enables better decisions. More pointedly: \textbf{uncertainty
is not a property of a model, but of a decision process.} Two models with
identical calibration can produce decisions of very different quality,
because the value of an uncertainty map depends on the rule that consumes
it. The same uncertainty map can look mediocre under one deferral rule
and excellent under another. These are not calibration statements; they
are statements about the joint estimation--decision system, and no
amount of recalibration addresses them.

This paper operationalizes that reframing. We study, empirically, how the
choice of uncertainty source and the choice of deferral policy jointly
determine segmentation decision quality, and we show that ignoring either
dimension leaves most of the available error reduction on the table. The
right combination of uncertainty method and deferral policy eliminates
roughly 80\% of segmentation errors while deferring only 25\% of pixels;
the wrong combination reduces errors by less than a third of that, using
the same images, the same model, and the same forward passes. The
difference is entirely in the \emph{decision layer} sitting on top.

\paragraph{Why this is not obvious.} One might object that ``use
uncertainty to defer'' is folklore, and therefore our reframing merely
restates what the field already knows. We disagree. In practice, existing
work optimizes uncertainty \emph{in isolation}: Bayesian methods are tuned
for calibration, ensembles for NLL, post-hoc methods for ECE. Deferral,
when it appears at all, is a fixed-threshold afterthought evaluated at a
single operating point. This decomposition hides the interaction we study.
A well-calibrated model with a poor deferral rule can underperform a
less-calibrated model with a well-designed one; any methodology that
scores uncertainty on its own is blind to this fact. The standard
calibration toolkit is not merely incomplete as a proxy for decision
quality---it is actively misleading, because it rewards distributional
match over discriminative separation, and deferral cares only about the
latter. Our contribution is to measure the interaction directly,
demonstrate that it is large, and show that the field has been optimizing
the wrong objective.

\paragraph{Setup.} Our study uses retinal vessel segmentation on DRIVE,
with cross-dataset evaluation on STARE and CHASE\_DB1. We train a U-Net
with a ResNet-34 encoder and compare two uncertainty estimation
approaches---MC Dropout (30 stochastic forward passes) and Test-Time
Augmentation (6 geometric transforms)---under three deferral strategies
of increasing sophistication:

\begin{enumerate}[leftmargin=*]
  \item \textbf{Global thresholding}: defer all pixels whose uncertainty
  exceeds a fixed threshold, optimized on a validation set.
  \item \textbf{Image-adaptive thresholding}: compute a per-image threshold
  based on percentile statistics of that image's uncertainty distribution,
  adapting the deferral rate to each image's difficulty.
  \item \textbf{Confidence-aware deferral}: weight uncertainty by inverse
  prediction confidence, so that uncertain pixels near the decision
  boundary ($p \approx 0.5$) are deferred preferentially over uncertain
  pixels where the model is nonetheless confident.
\end{enumerate}

We evaluate these combinations through risk-coverage curves, deferral
curves, and direct measurement of error reduction as a function of
deferral rate.

\paragraph{Contributions.}
\begin{enumerate}[leftmargin=*]
  \item \textbf{An estimation$\,\to\,$decision reframing of uncertainty
  for segmentation.} We treat uncertainty estimation and deferral as a
  single pipeline and show their interaction is large enough that
  optimizing either alone is insufficient. Under this lens, uncertainty
  quality is not a scalar property of a model but a property of the
  decision rule it is paired with.

  \item \textbf{TTA dominates MC Dropout as a decision-time uncertainty
  source.} TTA uncertainty discriminates correct from incorrect pixels
  substantially better (Unc-AUROC 0.881 vs.\ 0.722) while running
  $3.1\times$ faster---establishing it as the preferred uncertainty source
  for deferral-based segmentation pipelines.

  \item \textbf{A confidence-aware deferral rule.} We introduce the score
  $s = u \cdot (1 - c)$, which suppresses deferral for pixels that are
  uncertain but not borderline. Paired with TTA, it removes 55\% of errors
  at only 12\% deferral---the best low-budget efficiency in our study. TTA
  with image-adaptive deferral reaches roughly 80\% error reduction at
  25\% deferral---the overall error-removal ceiling.

  \item \textbf{Evidence that calibration is the wrong target for
  deferral.} Temperature scaling does not improve deferral operating
  points and can worsen ECE, demonstrating that calibration metrics and
  decision quality are decoupled and must be evaluated separately.

  \item \textbf{Cross-dataset evidence that decision quality is robust to
  domain shift.} Zero-shot evaluation on STARE and CHASE\_DB1 shows
  Unc-AUROC holding above 0.83 despite a 6.5-point Dice drop. The decision
  layer remains reliable exactly when the estimation layer is most stressed.
\end{enumerate}

The remainder of the paper proceeds as follows.
Sections~\ref{sec:related_work}--\ref{sec:experimental_setup} situate the
work, formalize the two-stage pipeline, and describe the experimental
protocol. Sections~\ref{sec:results}--\ref{sec:ablations} present results
and ablations. Sections~\ref{sec:discussion}--\ref{sec:conclusion} discuss
implications, limitations, and concluding arguments.
